%% ============================================================
%%  Chapter 1 — Introduction
%%
%%  Author      : Carlos David Prado-Socorro
%%  Date        : April 2026
%%  Institution : ICMol, Universitat de València
%%
%%  Compile with: pdflatex → biber → pdflatex × 2
%%  (or use latexmk -pdf -biber)
%% ============================================================

%% ---- If using as standalone file ----
\documentclass[12pt, a4paper, twoside]{report}

\usepackage[T1]{fontenc}
\usepackage[utf8]{inputenc}
\usepackage{lmodern}
\usepackage{microtype}

%% Page geometry
\usepackage[
  top=2.5cm, bottom=2.5cm,
  inner=3.0cm, outer=2.5cm,
  headheight=15pt
]{geometry}

%% Mathematics
\usepackage{amsmath, amssymb, amsthm}
\usepackage{bm}
\usepackage{siunitx}
\sisetup{
  separate-uncertainty = true,
  range-phrase = \text{--},
  detect-all
}

%% Chemistry (included for consistency with later chapters; not used in Part 1)
\usepackage[version=4]{mhchem}

%% Figures and tables
\usepackage{graphicx}
\usepackage{float}
\usepackage{booktabs}
\usepackage{array}
\usepackage{tabularx}
\usepackage{multirow}
\usepackage{caption}
\usepackage{subcaption}
\captionsetup{font=small, labelfont=bf, width=0.9\linewidth}

%% Colours
\usepackage[dvipsnames]{xcolor}

%% References and links
\usepackage[
  backend=biber,
  style=nature,
  sorting=none,
  maxnames=6,
  minnames=3,
  doi=true,
  url=false
]{biblatex}
\addbibresource{../bibliography/references.bib}

\usepackage[hidelinks, colorlinks=false]{hyperref}
\usepackage{cleveref}

%% Headers and footers
\usepackage{fancyhdr}
\pagestyle{fancy}
\fancyhf{}
\fancyhead[LE]{\small\leftmark}
\fancyhead[RO]{\small\rightmark}
\fancyfoot[C]{\small\thepage}

%% Paragraph spacing
\setlength{\parindent}{0pt}
\setlength{\parskip}{6pt}

%% ============================================================
\begin{document}

%% ============================================================
\chapter{Introduction}\label{ch:intro}
%% ============================================================

The research problem addressed in this thesis is the design, fabrication, and characterisation of a class of electronic devices --- organic memristive devices based on ion-containing polymer composites --- whose dynamical behaviour can be exploited as a computational resource. Because these devices occupy an unusual position at the intersection of solid-state electronics, polymer physical chemistry, and computational neuroscience, no single disciplinary perspective is sufficient to establish the context in which the rest of the thesis operates. The present chapter is therefore arranged as a sequence of complementary surveys. The remainder of this chapter addresses, in turn, the structural limitations of contemporary computing architectures and the motivation to move beyond them (\cref{sec:crisis}); the biological analogy from which neuromorphic engineering has historically drawn inspiration; the physics and taxonomy of memristive devices as candidate electronic primitives for post-von-Neumann computing; and finally the specific class of organic polymer-composite devices studied in this thesis, together with the explicit scientific objectives that the remaining chapters address. The statement of thesis objectives in the closing section fixes the scope of the experimental work and defines the criteria by which the contribution of this thesis is to be evaluated.


%% ------------------------------------------------------------
\section{The Crisis in Modern Computing}\label{sec:crisis}
%% ------------------------------------------------------------

The objective of the present section is to set out, in terms accessible to a reader trained primarily in physical chemistry rather than in computer architecture, why the hardware of modern digital computers is widely regarded as approaching a structural impasse, and why that impasse motivates the exploration of fundamentally different electronic primitives. Two lines of argument are pursued in parallel. The first, developed in \cref{subsec:vonneumann}, traces the limitations of the stored-program architecture that has served as the basis of essentially all general-purpose computing machines since the late 1940s, and documents their quantitative manifestation in the present era as the \textit{memory wall}, the \textit{energy wall}, and the end of classical transistor scaling. The second, developed in \cref{subsec:paradigms}, surveys two distinct responses to those limitations that have emerged in the device-physics community: the integration of computation and storage within a single physical substrate (in-memory and analogue computing), and the abandonment of clocked, synchronous arithmetic in favour of event-driven, temporally structured operation. This second response --- and in particular the identification of fading memory as a computational resource rather than an engineering defect --- supplies the conceptual frame within which the devices reported in this thesis are to be evaluated.

\subsection{Von Neumann Architecture and Its Limitations}\label{subsec:vonneumann}

The dominant architectural model underlying modern digital computers was articulated by von~Neumann in the ``First Draft of a Report on the EDVAC''~\cite{VonNeumann1993}, circulated in 1945 and formalised over the subsequent decade. Its defining feature is the \textit{stored-program} concept: instructions and data are represented in the same memory and drawn from it, one word at a time, by a central processing unit under the control of an instruction pointer. This unification of code and data was not an incidental engineering choice but the enabling abstraction of general-purpose computing: a single physical machine becomes capable of implementing any computable function by loading a different sequence of binary words into memory. The resulting architecture --- comprising a central processing unit, a main memory holding both code and data, and a bus connecting them --- is so deeply embedded in the design of contemporary digital systems that it is usually referred to simply as \textit{the} architecture of computing, and its constituent elements (cores, caches, dynamic random-access memory, instruction fetch, register file) are treated as unquestioned primitives in the design of processors for server, desktop, mobile, and embedded applications. More than seventy years of cumulative engineering effort have been expended on optimising every stage of this pipeline while leaving its basic topology essentially untouched.

A direct consequence of the stored-program model is that the processor and the memory are \textit{physically distinct} subsystems, separated by a communication channel whose bandwidth is a primary determinant of overall system performance. Every elementary computational step --- the fetch of an instruction, the load of an operand, the write-back of a result --- requires one or more traversals of this channel. In a modern general-purpose processor, the channel is a hierarchy of caches bridging between register files clocked at several gigahertz and main memory implemented in dynamic random-access memory (DRAM), which operates at effective bandwidths one or two orders of magnitude lower. Even when the entire working set of a computation fits into cache, the separation of compute and storage imposes an unavoidable cost: data must be moved, in the form of electrical signals along interconnect metal, between the transistors that perform arithmetic and the capacitors or latches that hold state. No arithmetic operation can be performed without first transporting operands, and no result can be preserved without transporting it back. This fundamental pattern is usually referred to in the architecture literature as the \textit{processor--memory dichotomy}.

In its simplest form, the von~Neumann architecture exposes only a single port between processor and memory, so that at any given clock cycle a single word can be read or written but not both simultaneously. The resulting choke point is the \textit{von~Neumann bottleneck}, a term introduced by Backus in the late 1970s and used throughout the subsequent literature to describe the limitation that even an arbitrarily fast processor can make no more progress than the memory interface permits~\cite{NatureBottleneck2018}. Decades of architectural innovation --- pipelined instruction fetch, multi-level cache hierarchies, speculative execution, prefetching, out-of-order issue, banked and multi-ported memory subsystems, and dedicated coprocessors ranging from vector units to graphics processors --- can be read as increasingly elaborate attempts to mitigate this bottleneck without replacing the underlying topology. Each such innovation preserves the fundamental distinction between compute and storage while partially hiding it from the software layer by amortising data movement over larger blocks, deeper queues, or more parallel channels. None removes the underlying asymmetry, and each carries its own overhead in silicon area, power, and design complexity. The cumulative effect is that the bottleneck has been displaced and renamed rather than dissolved; it reappears, in each successive architectural generation, at whichever interface happens to be slowest and most energy-expensive at that point in time.

The quantitative severity of the processor--memory imbalance was diagnosed explicitly by Wulf and McKee~\cite{WulfMcKee1995} in 1995, in an analysis that has come to be known as the \textit{memory wall}. They observed that processor performance, as measured by the product of clock frequency and instructions per cycle, had been improving at an annual rate of approximately 60\%, whereas the access latency and bandwidth of commodity dynamic memory had improved by close to an order of magnitude less per year. A simple arithmetic extrapolation of the two curves, taken from any plausible starting year, leads within a decade or two to a regime in which the average cycle is dominated almost entirely by memory access: the processor spends far more time waiting for operands than it does computing with them. Empirical evidence accumulated over the following two decades has confirmed this projection in detail. Modern microprocessors typically allocate between 30\% and 60\% of their total silicon area, and an even larger fraction of their power budget, to cache memory, prefetch logic, memory controllers, and interconnect, precisely because those resources are required to keep the arithmetic units supplied with data. Even in heavily cache-optimised numerical codes, a non-negligible fraction of cycles is spent stalled on memory operations, and the situation is considerably worse for workloads whose access pattern is irregular or whose working set exceeds the on-chip cache. The memory wall is thus not a speculative prediction but a measurable property of present-day systems, and one that scales unfavourably with computation size.

The memory wall as originally formulated was a performance argument. A second, and in several respects more fundamental, argument concerns energy. In a quantitative survey presented as a plenary address at the 2014 International Solid-State Circuits Conference, Horowitz summarised the cost of elementary operations in a representative modern CMOS process~\cite{Horowitz2014}. A $32$-bit integer addition consumed of the order of \SI{0.1}{\pico\joule}; a $32$-bit floating-point multiplication, a few picojoules; a $32$-bit register-file read, also a few picojoules. By contrast, fetching a $32$-bit word from on-chip static random-access memory cost tens of picojoules, and fetching the same word from off-chip DRAM cost several hundred picojoules to nanojoules --- that is, three to four orders of magnitude more energy than the arithmetic operation that the word is fetched to enable. The consequence is that, in any workload dominated by data movement, the total energy budget is set not by the cost of computing results but by the cost of shuffling operands. This asymmetry is the \textit{energy wall}, and it is not a transient property of a particular process node. It reflects an underlying physical fact: moving charge across a wire of non-trivial length dissipates energy in the wire's parasitic capacitance, and the total wire length in a memory subsystem greatly exceeds that within an arithmetic unit. The implication is direct: any architecture in which computation and storage are spatially separated is fundamentally bounded, in both speed and energy, by the cost of moving bits between them, and that cost does not scale away with transistor miniaturisation.

For several decades, the architectural inefficiencies catalogued above were absorbed, rather than eliminated, by the relentless exponential improvement of the underlying silicon technology. Two empirical scaling laws describe the regime in which that absorption was possible. The first, associated with Moore~\cite{Moore1965}, observed that the number of transistors that can be economically integrated on a single chip doubles at roughly constant intervals of twelve to twenty-four months. The second, due to Dennard and co-workers~\cite{Dennard1974}, provided the physical rationale for why that doubling should translate into proportionate performance gains: if every linear dimension of a MOSFET is scaled by a factor \(1/k\) and the supply voltage is scaled by the same factor, then the transistor switches \(k\) times faster, consumes \(1/k^{2}\) the energy per switching event, and requires \(1/k^{2}\) the area, so that the power density of a circuit at constant activity remains invariant. When the two laws held simultaneously, each silicon generation delivered both more transistors and faster, cooler ones, and the delivered performance of a given microarchitecture could be improved without structural change. The memory wall and energy wall were then, in practice, perceived as distant rather than binding limits, because each generation of processor was fast enough to make its own memory subsystem look slower than the previous generation's, without that contrast proving fatal for the workloads of the day.

Dennard scaling, however, is known to have broken down progressively from approximately 2005 onwards, when the gate-oxide thickness approached a regime in which quantum-mechanical leakage currents dominate sub-threshold behaviour and prevent further reduction of the supply voltage at the rate required by the scaling theory~\cite{Hennessy2019}. Transistor count has continued to grow, but the associated performance and energy-efficiency gains per transistor have slowed markedly, and the power density of a uniformly active die has begun to increase with each generation --- the regime commonly described as \textit{post-Dennard} or, with an allusion to the resulting necessity of selectively deactivating portions of the die to remain within a thermal budget, the \textit{dark silicon} era. Under these conditions, the previously tolerated inefficiencies of the von~Neumann architecture become first-order constraints. The hardware community can no longer rely on a new process node to mask the cost of data movement; that cost must instead be addressed architecturally. Simultaneously, and essentially independently, the workloads that dominate modern computing have shifted away from the single-thread, cache-friendly scientific and business codes for which the classical memory hierarchy was optimised, and towards machine-learning training and inference kernels that are defined by enormous repeated matrix multiplications over parameter tensors resident in memory~\cite{LeCun2015}. These workloads intensify, rather than relax, the data-movement bottleneck, and they do so at a historical moment at which the underlying silicon can no longer be relied upon to compensate.

The combination of these three trends --- the structural dominance of data movement over arithmetic, the corresponding dominance of memory-access energy over arithmetic energy, and the end of the Dennardian compensation mechanism --- is what justifies describing modern computing as facing a \textit{structural} rather than an \textit{incremental} problem. It is increasingly difficult to see further progress as a matter of refining the von~Neumann pipeline; the architecture itself, and in particular its physical separation of computation from storage, is now understood as the central obstacle. This recognition is not a minority view: it has been articulated repeatedly, including in the Turing Award lecture of Hennessy and Patterson~\cite{Hennessy2019}, in which the authors argue explicitly that the combination of the end of Dennard scaling and the rise of machine-learning workloads represents an inflection point comparable to the transition from vacuum tubes to transistors, and that the next decade of computer architecture will be defined by domain-specific hardware that co-locates memory and computation. It is against this background that the search for new electronic primitives --- devices whose operation does not presuppose the traditional distinction between storage and processing, and whose physics is naturally suited to supporting large, low-energy data-dependent operations --- has moved from an exotic research programme to a mainstream concern.

\subsection{In-Memory and Event-Driven Computing Paradigms}\label{subsec:paradigms}

Two broad responses to the situation described in \cref{subsec:vonneumann} have emerged within the device-physics and circuits communities over the past decade. Both abandon, in their most consistent formulations, the strict separation of computation and storage that defines the classical von~Neumann machine, but they do so along different axes and with different computational targets. The first response retains the clocked, arithmetic character of digital computing but dissolves the spatial separation between memory and processor, performing arithmetic operations directly within (or immediately adjacent to) the array elements that hold the operands; this family of approaches is known collectively as \textit{in-memory}, \textit{near-memory}, or \textit{processing-in-memory} computing. The second response retains, or even deepens, the spatial separation between different functional elements but abandons the synchronous, arithmetic model of computation in favour of operation driven by temporally structured, typically sparse, event streams, and exploits the intrinsic dynamical behaviour of physical devices as part of the computation itself. This second family is known as \textit{event-driven} or \textit{temporal} computing.

The central conceptual step of in-memory computing is the replacement of the classical \textit{fetch--compute--write-back} loop by an operation in which the computation is implemented \textit{physically} within a memory array. The canonical example is the analogue matrix--vector multiplication performed on a resistive crossbar. In such an array, a two-dimensional grid of memory cells holds conductance values that encode the entries of a matrix; the input vector is presented as a pattern of voltages applied to the rows, and the output vector appears as the column currents dictated by Kirchhoff's laws, each column current being the sum of row voltages multiplied by the corresponding conductances. A single clock period, and an energy cost dominated by the steady-state dissipation in the array rather than by any per-operation data transfer, thus yields a result that on a conventional processor would require as many multiply--accumulate operations as the matrix has entries. Demonstrations of this principle in hardware based on metal-oxide resistive memories~\cite{Prezioso2015}, phase-change materials~\cite{Koelmans2015, Sebastian2014}, and organic ionic devices~\cite{Fuller2019, VanDeBurgt2017, VanDeBurgt2018} have shown that matrix--vector products of practical size can be performed at energy efficiencies far beyond what digital accelerators based on standard CMOS logic can achieve. The same mechanism can be generalised to support convolutional operations, attention mechanisms, and other primitives central to contemporary machine learning. In this framing, the memory array is no longer a passive store but an active computational substrate, and the requirements placed on its individual cells are correspondingly different from those placed on the cells of a classical memory.

The device-level requirements imposed by the in-memory-computing paradigm, in its standard form, nevertheless remain recognisably those of a memory. A single cell should store a quasi-static conductance value that encodes a trained weight; that value should be set once (or updated occasionally during training) and then read, potentially many millions of times, without drifting significantly. The cell should be reproducible across a large array, with device-to-device and cycle-to-cycle variability kept below a tolerance set by the numerical precision demanded by the application. Non-volatility over operational timescales of seconds to years is a design goal: a cell that loses its stored value within milliseconds is, under these criteria, defective. When organic memristive devices are evaluated against this yardstick, the comparison is frequently unflattering. Polymer-composite ionic devices tend to show substantial device-to-device variability, conductance drift during repeated reads, and retention times of seconds to hours rather than years. An evaluation that takes the in-memory, non-volatile weight-storage criterion as its point of reference is therefore likely to classify such devices as promising but ultimately flawed competitors to metal-oxide or phase-change memories. It is important to recognise, however, that this is the evaluation criterion of a \textit{specific} computational paradigm, and that alternative paradigms exist within which the same physical properties --- short retention, device-to-device heterogeneity, dynamical relaxation --- are not liabilities but resources.

The second response to the limitations of the von~Neumann model, often developed in parallel to the first, begins from the observation that not all useful computation is most naturally expressed as a sequence of arithmetic operations on static operands. Many tasks of interest --- the processing of natural auditory and visual signals, the control of continuous motor systems, the classification of biomedical waveforms, the extraction of features from streaming sensor data --- are defined on inputs that are themselves time-varying and that carry information in their temporal structure. For such tasks, a natural computational model is one in which the inputs are represented as streams of \textit{events} in continuous time, each event consisting of a channel identity and a timestamp, and in which the state of the computing system evolves continuously in response to these events rather than stepping through discrete clock cycles. This model is realised at the algorithmic level by spiking neural networks, at the circuit level by event-driven asynchronous digital logic, and at the device level by physical substrates whose intrinsic relaxation dynamics are of the same order of magnitude as the temporal structure of the inputs they process. It is this last level on which the present thesis operates.

The loose inspiration for this second paradigm is, characteristically, biological. Mammalian nervous systems perform demanding sensorimotor and cognitive tasks at power budgets several orders of magnitude below those of the most efficient artificial counterparts, and it is difficult to avoid the conclusion that at least part of this efficiency is attributable to architectural choices that differ systematically from the von~Neumann model. Biological neurons emit discrete events (action potentials) sparsely in time; communicate through synapses whose conductance is itself a dynamical variable rather than a static weight; maintain and modify their state through slow, graded electrochemical processes rather than through explicit write operations; and are embedded in a physical medium (the extracellular ionic environment) whose properties are part of the computation. Although a direct transcription of these mechanisms into silicon is neither feasible nor obviously desirable, several architectural lessons are widely acknowledged as transferable: representing information as sparse events rather than dense words~\cite{IndiveriLiu2015}, exploiting slow device physics as a computational primitive rather than engineering it away, and treating device-to-device heterogeneity as a source of useful functional diversity rather than as noise to be suppressed. These lessons define the conceptual basis of neuromorphic engineering as a discipline, and they constitute the second route out of the von~Neumann crisis pursued in the device-physics literature.

A particularly clean theoretical articulation of the computational value of dynamical devices was provided by Boyd and Chua~\cite{BoydChua1985}, who showed that any time-invariant nonlinear operator with \textit{fading memory} on a suitable function space can be uniformly approximated by a finite Volterra series. The statement carries a direct engineering corollary: a physical system whose present state carries a smoothly decaying trace of its past inputs can, when combined with a trainable linear readout, implement a broad class of useful time-varying computations. This insight was later operationalised, independently, by Jaeger in the form of the \textit{echo state network}~\cite{Jaeger2001ESN, JaegerHaas2004} and by Maass, Natschläger and Markram in the form of the \textit{liquid state machine}~\cite{Maass2002LSM}, and subsequently unified under the name \textit{reservoir computing}~\cite{LukoseviciusJaeger2009, Verstraeten2007}. In the reservoir-computing paradigm, an untrained dynamical system --- the reservoir --- is driven by the input stream, and only a linear readout layer on top of its state is adjusted by learning. When the reservoir is realised in physical hardware rather than as a simulated recurrent network, the paradigm is called \textit{physical reservoir computing}~\cite{Tanaka2019}, and its hardware instances have included optical, mechanical, and spintronic systems as well as electronic devices whose internal dynamics exhibit the necessary fading-memory property. Demonstrations of physical reservoir computing in volatile memristive~\cite{Du2017, Midya2019, Zhong2021} and in organic electrochemical~\cite{Cucchi2021} substrates are of direct relevance to the device class studied in this thesis, and will be revisited when the computational framework is developed in a later chapter.

The central implication of this second paradigm for the present thesis is a reframing of the criteria by which a candidate memristive device is to be judged. Under the in-memory, non-volatile-weight paradigm, a short retention time is a defect: the device loses information that the application requires it to keep. Under the physical-reservoir and event-driven-computing paradigm, by contrast, a finite retention time is the very property that makes the device computationally useful, because it defines the effective window over which past inputs influence the present output --- that is, the fading-memory timescale. Similarly, device-to-device variability, which is penalised under the weight-storage paradigm, becomes a resource under the reservoir paradigm: a bank of devices with a spread of intrinsic timescales, coupled to a common linear readout, can jointly resolve temporal structure over a wider range of scales than any single device could~\cite{Tanaka2019}. The same physical property --- for example, the ionic relaxation time of a polymer electrolyte --- that makes a device a poor non-volatile memory can therefore make it an excellent temporal computing primitive, provided that the architecture in which it is embedded is one that exploits, rather than suppresses, its dynamics. The work reported in this thesis adopts this second viewpoint throughout: the devices studied are not candidate non-volatile weights, and they are not evaluated as such. They are candidate heterogeneous, fading-memory elements, and the criteria by which they succeed or fail are those appropriate to that class.

Two broad commitments therefore shape the remainder of this chapter and the experimental programme to which it serves as an introduction. The first is that the computational target is not non-volatile weight storage within an analogue crossbar but temporal computing within an event-driven or reservoir-style architecture; the second is that the evaluation criteria applied to the devices studied here are those appropriate to temporal computing, and that comparisons to fast, high-endurance, non-volatile memories are set aside as category errors rather than as fair contests. The scientific contribution documented in the later chapters is, accordingly, structured in three layers: a single organic memristive device that is instrumented deeply enough to establish that the physical mechanism invoked in its design is in fact at work (reported in \cref{ch:poc}); a broader comparative corpus of polymer-electrolyte devices, studied across composition and cation identity on a restricted but internally consistent set of dynamical measurements, that quantifies how the resulting fading-memory behaviour can be tuned by chemistry; and a data-driven demonstration that those chemistry-tuned dynamics can be exploited, in simulation on the measured data, as primitives of temporal computation. The remaining sections of the present chapter build up the physical and biological context required for this programme, and close with an explicit statement of the objectives against which the thesis is to be assessed.

\end{document}
